\paragraph{Data sources and purpose.}
Our main source is Wright and Augenstein's abstract--press-release dataset~\cite{wright2021}.
Each pair contains a press-release claim, an aligned abstract excerpt, two language-strength codes, and an exaggeration label.
These labels compare wording strength: stronger, equal, or weaker than the source.
They support detection experiments, but do not provide evidence passages or reference corrections for evaluating edits.

\begin{center}
\begin{tabular}{lrrrr}
\toprule
Original split & Exaggerates & Same & Downplays & Total \\
\midrule
Training & 22 & 61 & 17 & 100 \\
Test & 122 & 345 & 96 & 563 \\
\midrule
Total & 144 & 406 & 113 & 663 \\
\bottomrule
\end{tabular}
\refstepcounter{table}\label{tab:paired-data}
\par\smallskip
Table \thetable: Original pairs. The prepared sets below are subsets, not additional pairs.
\end{center}

Equal strength does not guarantee support: the population or outcome may still be overstated.
We therefore add separate editing judgments while retaining every original row and label.

The auxiliary collection contains 3,061 abstract sentences and 2,076 press-release sentences~\cite{yu2019,yu2020}.
These 5,137 individual sentences can support language-strength learning.
They cannot test whether a claim matches a particular source.
Their categories are non-claim, association, conditional causation, and direct causation.
Auxiliary label 0 means non-claim; paired-data strength 0 means an explicit statement of no relationship.
We keep these meanings separate.

\paragraph{Development examples and annotation rules.}
We prepared eight training examples: four originally labeled exaggerates and four labeled same.
Eight follows the team's five-to-eight-example review plan.
Selection covers decisions, not population error rates.
We recovered complete abstracts through Europe PMC.
Records include publication identifiers, source versions, exact evidence passages and offsets, working decisions, proposed corrections, and reasons.

We assign three decisions: supported, overclaim, and insufficient evidence.
Supported claims should remain exactly unchanged.
Overclaims need the smallest supported correction.
If evidence is insufficient, the editor should explain why and withhold a correction.
Reviewers check quantities, population, outcome, comparisons, and uncertainty alongside causal wording.
For example, one working revision replaces ``is saving lives'' with ``was associated with lower mortality.''
Other cases require preserving justified stronger wording or recognizing missing evidence.
Full-abstract editing decisions can differ from original excerpt-based labels.
The current decisions are three overclaims, three supported claims, and two insufficient-evidence cases.
Independent human review is pending.
All eight remain development-only; their counts are not model results.

\paragraph{Pilot selection, review, and remaining work.}
The 44 pilot candidates comprise all 22 training exaggerates cases and 22 same cases.
Supported candidates test unnecessary editing.
Controls include four development examples and 18 cases selected by fixed hash ranking.
All eight development examples are already included in the 44.
Original-label balance may change after review.
Scope-only errors and insufficient-evidence cases will be analyzed separately from causal errors.

Eight candidates have full abstracts; source recovery remains for 36.
Two people will independently judge each usable pair and verify its evidence and correction.
We will retain initial judgments and disagreements.
We plan about 20 development cases and up to 24 evaluation cases after recovery and review.
The split is not frozen.
Each verified paper must belong to one split.
Papers already used in development remain there.
The 563 original test pairs are reserved from model and prompt development.
Their labels support detection testing, not editing evaluation.
Detector inputs contain only claim and source text; editor inputs contain claim and assigned evidence.
Reference decisions, strength codes, and corrections are never prediction inputs.
Before training, we will check auxiliary data for overlap with reserved papers and texts.

The proposal targeted 150--250 newly annotated examples.
Dataset preservation, eight development records, and the 44-case acquisition pool are complete.
A human-verified editing evaluation set remains unfinished.
Next come source recovery, independent review, and frozen paper-separated splits before experiments.
Expansion toward 150--250 depends on collection and annotation capacity after this pilot.
